\section{版本、目标插值与重复协议核对}
\label{app:alignment}

\paragraph{三个不同层面的实验。}
架构移植保留 \method{} 的共同迭代流程，并按模型修改教师输出及可合并层；FeatCal 移植则是本地 T5 基线。目标特征插值改变 \method{} 的拟合目标，是单独扩展。多种子固定方法与超参数、改变校准抽样，是重复协议。默认 \method{} 在 CLIP、T5 与 LLM 中均关闭目标插值；不能将 T5 的扩展结果默认为跨架构主方法。

\begin{table}[htbp]
\centering\small
\caption{当前实际覆盖。种子指校准抽样，不是评测器内部的随机数设置。}
\label{tab:protocol-coverage}
\setlength{\tabcolsep}{4pt}
\begin{tabular}{@{}lll@{}}
\toprule
实验组 & 目标插值 & 已完成种子 \\
\midrule
CLIP 工作台九格 & 关闭 & 每格只有 0 \\
正式 CLIP B/32 八任务 & 关闭 & 0、1、2、3 \\
T5-base 默认 \method{} / FeatCal & \method{} 关闭 & 两者各 0、1、2、3、4 \\
T5-base \method{} 插值 0.3 / 0.6 & 开启 & 分别 0--3 / 0--2 \\
T5-large 默认 \method{} / FeatCal & \method{} 关闭 & 两者各只有 0 \\
T5-large \method{} 插值 0.3 / 0.6 & 开启 & 0.3 只有 0；0.6 缺完整评测 \\
LLM 十组起点 & 关闭 & 每组只有 42 \\
Llama Iso-C 校准重复 & 关闭 & 42、1、2，仅 GSM8K 齐全 \\
\bottomrule
\end{tabular}
\end{table}

% Generated by motivation_method/update_results.py from data/verified_results.json.
\begin{table}[htbp]
\centering\small
\caption{T5-base 目标插值消融：三行统一使用已完成的配对 seed 0、1、2。其它 \method{} 配置相同；误差为样本标准差。}
\label{tab:t5-target-paired}
\begin{tabular}{@{}lcc@{}}
\toprule
\method{} 目标版本 & seeds & 八任务宏平均 \\
\midrule
原始目标 & 0, 1, 2 & $84.89\pm 0.04$ \\
目标插值 0.3 & 0, 1, 2 & $85.64\pm 0.05$ \\
目标插值 0.6 & 0, 1, 2 & $85.77\pm 0.08$ \\
\bottomrule
\end{tabular}
\par\smallskip
{\footnotesize 共同配置：Iso-C 1.3 起点，每任务统计 256、独立留出 256，四轮预算，CG 100 步，留出教师 KL 选步与选轮。完整 5/4/3 次记录保留在表~\ref{tab:t5-main}，不混成同一重复集合；两种扩展尚未补齐五种子。}
\end{table}


\paragraph{比较边界。}
T5-base 默认 \method{} 与 FeatCal 已有五个同编号种子，均值只差约 $0.079$ 个百分点，分别赢得三个、两个种子。两者起点及实际使用的选择数据量不同，因此该结果属于默认流程比较。检查 base 五种子和 large 单种子的八任务缓存，共 48 组，统计样本的索引与 token 均一致，\method{} 留出样本与统计样本互斥。不能由此把默认 FeatCal 描述成也使用了额外留出选择。

CG 次数相同也不表示求解协议完全相同。正式 CLIP/T5 以内层当前权重为初值；工作台/LLM 历史实现以内层专家平均权重为初值。外层 Iso-C 起点不能替代这项区分。CLIP 的逐块搜索与语言模型的全局搜索、教师构造及输出头处理也分别登记；不同种子选中不同轮数是同一选择规则的输出，不属于换超参数。

\paragraph{后续对齐规范与缺口。}
后续统一种子集合为 0--4，默认关闭目标插值，固定 Iso-C 1.3、统计与留出各 256、四轮预算、CG 100 步，并统一从当前权重开始内层求解。T5 目标消融单列原始目标、0.3、0.6；其它样本量、轮数与初始化干预只改变所研究因素。同设定内固定模型与数据指纹、样本 ID、教师、评测器、指标及选择规则。基线保留各自方法定义，同时披露数据和选择预算；等起点、等总数据预算的比较另设控制。

这是审计后制定的执行规范，不是历史预注册或已完成实验。T5 扩展仍缺种子，large 大多只有单种子；工作台与 LLM 修改 CG 初值后需要重跑完整重复集合。已有 seed 42 的 LLM 结果保留原编号。确定性基线不通过重复复制制造误差条，缺测也不进入均值。完整运行清单、配置来源与待执行项随稿件的协议审计文件提供。
